% ============================================================
%  Поглавље 7: Веза са аутоматским резоновањем
% ============================================================
\section{Квантно рачунарство и аутоматско резоновање}
\label{sec:ar}

% --- Циљ поглавља ---
% Повезати квантне технике са темом семинара: аутоматским резоновањем.
% Показати да постоје дубљи концептуални мостови, а не само
% аналогија по примени.

\subsection{SAT и квантно решавање}

Гроверов алгоритам претражује неструктурисани простор $N$ стања
у $O(\sqrt{N})$ корака, наспрам класичних $O(N)$~\cite{grover1997}.
За SAT са $n$ варијабли то даје убрзање:
\[
  O(2^n) \;\longrightarrow\; O(2^{n/2}).
\]

Интуиција је следећа. Булова формула
\[
  \varphi(x_1,\dots,x_n)
\]
посматра се као функција која за сваку доделу
$\mathbf{x}\in\{0,1\}^n$ враћа $1$ ако је формула задовољена,
а $0$ иначе. Тада се конструише \emph{квантни оракл}
\[
  O_\varphi \ket{\mathbf{x}} = (-1)^{\varphi(\mathbf{x})}\ket{\mathbf{x}},
\]
који не мења битски низ, већ мења \emph{фазу} оних стања која
представљају задовољавајуће доделе. Након тога Гроверова
итерација појачава амплитуде „добрих" стања и потискује остала,
па мерење са великом вероватноћом враћа решење.

Ако постоји $M$ задовољавајућих додела, сложеност је
$O(\sqrt{2^n/M})$, па је најчешће цитирани случај
$M=1$ или мало $M$, где добијамо $O(2^{n/2})$.
Међутим, ово \emph{не} значи да SAT постаје лак у полиномском смислу:
убрзање је квадратно, али је и даље експоненцијално, што је у складу
са релативизованим доњим границама за неструктурисану претрагу~\cite{bennett1997}.


\subsection{Хибридни квантно-класични приступи}

% --- Хибридни приступи ---
% Модеран приступ: QAOA и VQE (Variational Quantum Eigensolver) нису
% чисто квантни --- они користе класичне оптимизаторе за подешавање
% квантних параметара. Ово је директна аналогија са SMT решавачима
% које воде класични CDCL алгоритми.

% TODO: Нацртати дијаграм хибридне петље:
%   Квантно коло → Мерење → Класична оптимизација → (итерација)
% Поредити са CDCL петљом у SAT решавачима.

\subsection{SMT теорије и квантни Хамилтонијани}

Свака SMT формула може се посматрати као скуп ограничења над
променљивама, док се QUBO/Изингова формулација посматра као
енергетска функција чији минимум репрезентује решење~\cite{barrett2009,lucas2014,glover2022qubo}.
Та паралела је суштинска:
\begin{itemize}[nosep]
  \item SMT ограничења $C_i(\mathbf{x})$ одговарају penalty члановима
        $P_i(\mathbf{x})$ у QUBO Хамилтонијану,
  \item модел SMT формуле одговара основном стању
        (стању минималне енергије),
  \item теоријски решавач у SMT-у игра улогу „симболичког физичара":
        формално проверава да ли је кодирање коректно пре него што
        га предамо квантном или хеуристичком оптимизатору.
\end{itemize}

У том смислу, квантно жарење не замењује формалне методе, већ их
допуњује: квантни уређај претражује простор стања, док SMT решавач
проверава да ли енергетска функција заиста репрезентује оригинални
комбинаторни проблем.

\begin{figure}[h]
\centering
\begin{tikzpicture}[
  node distance=1.4cm and 0.9cm,
  box/.style={
    draw,
    rounded corners=4pt,
    thick,
    align=center,
    text width=2.45cm,
    minimum width=2.45cm,
    minimum height=1.0cm,
    font=\small,
    fill=gray!8
  },
  classical/.style={box, fill=blue!8, draw=blue!50!black},
  qubo/.style={box, fill=orange!12, draw=orange!70!black},
  quantum/.style={box, fill=red!8, draw=red!65!black},
  verify/.style={box, fill=green!10, draw=green!45!black},
  flow/.style={->, thick}
]
  \node[classical] (orig) {Полазни\\проблем};
  \node[classical, right=of orig] (smt) {SMT\\модел};
  \node[qubo, right=of smt] (qubo) {QUBO /\\Изинг};
  \node[quantum, right=of qubo] (solver) {Квантни\\решавач};
  \node[verify, below=1.4cm of qubo] (check) {Z3: провера\\решења};

  \draw[flow] (orig) -- (smt);
  \draw[flow] (smt) -- (qubo);
  \draw[flow] (qubo) -- (solver);
  \draw[flow] (orig.south) |- (check.west);
  \draw[flow] (solver.south) |- node[pos=0.3, right, font=\scriptsize] {$\hat{\mathbf{x}}$} (check.east);
\end{tikzpicture}
\caption{Хибридни ток рада у коме се полазни ограничени проблем најпре
преводи у SMT и QUBO/Изингову форму, затим решава квантно или
хеуристички, а добијени кандидат се на крају проверава помоћу Z3.}
\label{fig:qubo-z3-pipeline}
\end{figure}

Слика~\ref{fig:qubo-z3-pipeline} илуструје централну тезу овог рада:
QUBO није замена за аутоматско резоновање, већ интерфејс између
симболичког описа проблема и квантне оптимизације, док Z3 остаје
формални механизам који затвара петљу и проверава коректност добијених
решења.

\subsection{Z3 као верификатор QUBO ограничења}

Кључна слабост QUBO формулације је у томе што су ограничења
уграђена \emph{индиректно}, преко penalty чланова~\cite{lucas2014,glover2022qubo}.
Ако је знак погрешан, ако недостаје један квадратни члан, или ако је
коефицијент казне сувише мали, минимум енергије више не мора да
одговара важећем решењу оригиналног проблема.
Z3 је зато користан не само као решавач, већ и као
\emph{верификатор исправности саме QUBO трансформације}~\cite{demoura2008,bjorner2014}.

Нека је оригинални проблем дат у облику
\[
  \min g(\mathbf{x})
  \qquad \text{уз услове} \qquad
  C_1(\mathbf{x}), \dots, C_m(\mathbf{x}),
\]
а QUBO апроксимација као
\[
  f(\mathbf{x}) = g(\mathbf{x}) + \sum_{j=1}^m \lambda_j P_j(\mathbf{x}),
  \qquad \mathbf{x} \in \{0,1\}^n.
\]
Тада Z3 може да послужи у најмање три различита корака верификације.

\textbf{1. Провера семантике penalty члана.}
За свако ограничење треба доказати да penalty терм заиста
разликује дозвољена и недозвољена стања:
\[
  C_j(\mathbf{x}) \iff P_j(\mathbf{x}) = 0,
  \qquad
  \neg C_j(\mathbf{x}) \Rightarrow P_j(\mathbf{x}) > 0.
\]
На пример, за ограничење „тачно један" $x_i + x_j = 1$,
penalty члан
\[
  P(\mathbf{x}) = \lambda\,(1 - x_i - x_j + 2x_ix_j)
\]
је исправан ако Z3 покаже да је формула
\[
  x_i,x_j \in \{0,1\}
  \;\land\;
  \bigl((x_i+x_j=1 \land P \neq 0)
  \lor
  (x_i+x_j\neq 1 \land P = 0)\bigr)
\]
незадовољива. На тај начин добијамо \emph{контрапример ако кодирање
није добро}, уместо да се ослањамо само на ручни алгебарски развој.

\begin{example}[Провера penalty члана помоћу Z3]
\begin{lstlisting}[style=python, caption={Z3 provera da penalty clan korektno kodira ogranicenje $x_i + x_j = 1$.}]
from z3 import *

xi, xj, lam = Ints('xi xj lam')

# Penalty term for "exactly one"
P = lam * (1 - xi - xj + 2 * xi * xj)

s = Solver()
s.add(Or(xi == 0, xi == 1))
s.add(Or(xj == 0, xj == 1))
s.add(lam > 0)

# Search for a counterexample:
# valid assignment with nonzero penalty, or
# invalid assignment with nonpositive penalty
s.add(Or(
    And(xi + xj == 1, P != 0),
    And(xi + xj != 1, P <= 0)
))

print(s.check())  # unsat
\end{lstlisting}
\end{example}

Резултат \texttt{unsat} значи да контрапример не постоји:
penalty члан је нула тачно на дозвољеним стањима, а строго позитиван
на недозвољеним.

\textbf{2. Провера да је казна довољно велика.}
Чак и ако је $P_j(\mathbf{x}) = 0$ тачно на дозвољеним стањима,
коефицијент $\lambda_j$ мора бити довољно велик да ниједно
недозвољено стање не постане енергетски повољније од исправног решења~\cite{lucas2014,glover2022qubo}.
За мале и средње инстанце Z3 може прво да нађе оптимално дозвољено
решење $\mathbf{x}^\star$, а затим да провери да ли постоји
контрапример облика
\[
  \exists\, \mathbf{x}\in\{0,1\}^n :
  \neg C(\mathbf{x}) \land f(\mathbf{x}) \leq f(\mathbf{x}^\star).
\]
Ако је ова формула \texttt{unsat}, онда ниједно кршење ограничења
не може да „пробије" легално решење у QUBO енергији.
То је формални доказ да је изабрана penalty скала довољна.

\begin{example}[Провера да је penalty коефицијент довољно велик]
\begin{lstlisting}[style=python, caption={Z3 provera da penalizacija eliminise nedozvoljena resenja.}]
from z3 import *

# Original constrained problem:
# maximize 3*x0 + 2*x1  subject to x0 + x1 <= 1
x0, x1 = Ints('x0 x1')

opt = Optimize()
opt.add(Or(x0 == 0, x0 == 1), Or(x1 == 0, x1 == 1))
opt.add(x0 + x1 <= 1)

reward = 3 * x0 + 2 * x1
h = opt.maximize(reward)

assert opt.check() == sat
m = opt.model()
best_reward = m.eval(reward).as_long()   # 3
best_energy = -best_reward               # -3

def violates_feasible_optimum(lam_value):
    y0, y1 = Ints('y0 y1')
    s = Solver()
    s.add(Or(y0 == 0, y0 == 1), Or(y1 == 0, y1 == 1))

    qubo = -(3 * y0 + 2 * y1) + lam_value * y0 * y1

    # Look for an infeasible assignment that is at least as good
    # as the best feasible solution
    s.add(y0 + y1 > 1)
    s.add(qubo <= best_energy)

    return s.check()

print(violates_feasible_optimum(2))  # sat
print(violates_feasible_optimum(3))  # unsat
\end{lstlisting}
\end{example}

Овде \texttt{sat} за $\lambda = 2$ значи да и даље постоји
недозвољено стање које достиже енергију најбољег легалног решења,
док \texttt{unsat} за $\lambda = 3$ формално потврђује да је
та казна довољна.

\textbf{3. Верификација резултата квантног решавача.}
Када QAOA или квантни анилер врате кандидатно решење
$\hat{\mathbf{x}}$~\cite{farhi2014,preskill2018}, Z3 може независно да провери:
\begin{itemize}[nosep]
  \item да ли $\hat{\mathbf{x}}$ задовољава оригинална ограничења,
  \item колика је вредност оригиналне циљне функције $g(\hat{\mathbf{x}})$,
  \item да ли постоји боље легално решење за исту малу инстанцу.
\end{itemize}
Тако Z3 постаје \emph{референтни оракл} за процену квалитета
квантних и хеуристичких метода: квантни алгоритам тражи ниску
енергију, а SMT верификатор потврђује да та ниска енергија има
жељено семантичко значење.

\begin{example}[Верификација кандидатног bitstring-а]
\begin{lstlisting}[style=python, caption={Z3 provera kandidatnog resenja dobijenog kvantnim ili heuristickim resavacem.}]
from z3 import *

# Candidate returned by an annealer or QAOA
x_hat = [0, 1, 0]
weights = [5, 4, 1]

# Original problem: maximize 5*x0 + 4*x1 + x2
# subject to exactly one variable being 1
x = [Int(f'x{i}') for i in range(3)]

# 1) Check whether x_hat is feasible
s = Solver()
for v in x:
    s.add(Or(v == 0, v == 1))
for i, bit in enumerate(x_hat):
    s.add(x[i] == bit)
s.add(Sum(x) == 1)

print(s.check())  # sat

# 2) Evaluate the returned solution
hat_value = sum(weights[i] * x_hat[i] for i in range(3))
print(hat_value)  # 4

# 3) Search for a better feasible solution
opt = Optimize()
for v in x:
    opt.add(Or(v == 0, v == 1))
opt.add(Sum(x) == 1)

objective = Sum([weights[i] * x[i] for i in range(3)])
h = opt.maximize(objective)

if opt.check() == sat:
    m = opt.model()
    best_x = [m.eval(x[i]).as_long() for i in range(3)]
    best_val = m.eval(objective).as_long()
    print(best_x)    # [1, 0, 0]
    print(best_val)  # 5
\end{lstlisting}
\end{example}

У овом примеру решење $\hat{\mathbf{x}}$ јесте легално, али није
оптимално: Z3 налази бољи дозвољени битски низ и тиме пружа егзактан
класични референтни резултат за поређење.

\begin{remark}
У пракси је ово један од најважнијих хибридних образаца:
QUBO се користи као формат за оптимизацију, а Z3 као формални слој
који доказује да су penalty чланови коректни, да су коефицијенти
довољни и да добијени битски низ заиста представља важеће решење
полазног проблема.
\end{remark}
